\chapter{Breuken en machten}\label{ch:g9:fractions}

\omterm{def:g9:fractions:fraction}{Breuken} en \omterm{def:g9:fractions:power}{machten} zijn het dagelijkse gereedschap van het rekenen: hoeveelheden
verdelen, verhoudingen vergelijken, heel grote of heel kleine getallen opschrijven.
Dit hoofdstuk zet de rekenregels ervoor op een rij --- met elke tussenstap
uitgeschreven --- en voert de \omterm{def:g9:fractions:scientific}{wetenschappelijke notatie} in.

\section{Rekenen met breuken}

\begin{definition}[Breuk]\label{def:g9:fractions:fraction}
Een \emph{breuk}\index{breuk} $\dfrac ab$ (met $b \neq 0$) is het \omterm{def:g3:division:remainder}{quotiënt} van
$a$ door $b$. Twee breuken zijn gelijk wanneer je de ene uit de andere krijgt door
\omterm{def:g4:fractions:def}{teller} en \omterm{def:g4:fractions:def}{noemer} met hetzelfde getal, verschillend van \omterm{def:g1:counting:zero}{nul}, te vermenigvuldigen
(of erdoor te delen):
\[
\frac ab = \frac{a \times k}{b \times k} \qquad (k \neq 0).
\]
\end{definition}

\begin{example}\label{ex:g9:fractions:simplify}
Vereenvoudig $\dfrac{42}{56}$ stap voor stap:
\[
\frac{42}{56} = \frac{42 \div 2}{56 \div 2} = \frac{21}{28}
= \frac{21 \div 7}{28 \div 7} = \frac34 .
\]
(In \cref{ch:g9:arith} doet de grootste gemene \omterm{def:g5:division:multiple}{deler} dit in één stap.)
\end{example}

\begin{method}[Breuken optellen en aftrekken]\label{met:g9:fractions:add}
\begin{enumerate}
  \item Zet beide \omterm{def:g9:fractions:fraction}{breuken} op een \emph{gemeenschappelijke \omterm{def:g4:fractions:def}{noemer}} (een
        gemeenschappelijk \omterm{def:g5:division:multiple}{veelvoud} van de twee \omterm{def:g4:fractions:def}{noemers});
  \item tel de \omterm{def:g4:fractions:def}{tellers} op of trek ze af, en houd de \omterm{def:g4:fractions:def}{noemer};
  \item vereenvoudig het resultaat als dat kan.
\end{enumerate}
\end{method}

\begin{example}\label{ex:g9:fractions:add}
Bereken $\dfrac56 + \dfrac34$. Een gemeenschappelijke \omterm{def:g4:fractions:def}{noemer} van $6$ en $4$ is
$12$:
\[
\frac56 + \frac34
= \frac{5 \times 2}{6 \times 2} + \frac{3 \times 3}{4 \times 3}
= \frac{10}{12} + \frac{9}{12}
= \frac{19}{12}.
\]
Bereken $2 - \dfrac37$. Schrijf $2$ als een \omterm{def:g9:fractions:fraction}{breuk} op \omterm{def:g4:fractions:def}{noemer} $7$:
\[
2 - \frac37 = \frac{14}{7} - \frac37 = \frac{11}{7}.
\]
\end{example}

\begin{omfigure}
\begin{tikzpicture}[scale=0.55]
  \foreach \i in {0,...,11}
    \draw[omExo] ({mod(\i,12)},0) rectangle ({mod(\i,12)+1},1);
  \foreach \i in {0,...,9}
    \fill[omDef!35] (\i,0) rectangle (\i+1,1);
  \draw[thick] (0,0) rectangle (12,1);
  \node[below, font=\small] at (6,-0.1)
    {$\frac56 = \frac{10}{12}$: 10 vakjes van 12};
  \begin{scope}[yshift=-2.4cm]
  \foreach \i in {0,...,11}
    \draw[omExo] (\i,0) rectangle (\i+1,1);
  \foreach \i in {0,...,8}
    \fill[omThm!35] (\i,0) rectangle (\i+1,1);
  \draw[thick] (0,0) rectangle (12,1);
  \node[below, font=\small] at (6,-0.1)
    {$\frac34 = \frac{9}{12}$: 9 vakjes van 12};
  \end{scope}
\end{tikzpicture}

{\small Waarom een gemeenschappelijke \omterm{def:g4:fractions:def}{noemer}: door beide gehelen in $12$ gelijke
vakjes te snijden worden de \omterm{def:g9:fractions:fraction}{breuken} vergelijkbaar en optelbaar --- samen
$10 + 9 = 19$ twaalfden.}
\end{omfigure}

\begin{method}[Breuken vermenigvuldigen en delen]\label{met:g9:fractions:mult}
\begin{enumerate}
  \item Om te vermenigvuldigen, vermenigvuldig je de \omterm{def:g4:fractions:def}{tellers} met elkaar en de
        \omterm{def:g4:fractions:def}{noemers} met elkaar: $\dfrac ab \times \dfrac cd = \dfrac{ac}{bd}$
        (vereenvoudig \emph{vóór} het vermenigvuldigen wanneer dat kan);
  \item om te delen, vermenigvuldig je met het omgekeerde van de \omterm{def:g5:division:multiple}{deler}:
        $\dfrac ab \div \dfrac cd = \dfrac ab \times \dfrac dc$.
\end{enumerate}
\end{method}

\begin{example}\label{ex:g9:fractions:mult}
\[
\frac{7}{15} \times \frac{25}{14}
= \frac{7 \times 25}{15 \times 14}
= \frac{7 \times 25}{15 \times 2 \times 7}
= \frac{25}{30} = \frac56,
\]
waar we eerst door de gemeenschappelijke factor $7$ en daarna door $5$ hebben
vereenvoudigd. En een deling:
\[
\frac{3}{4} \div \frac{9}{8}
= \frac34 \times \frac89
= \frac{3 \times 8}{4 \times 9}
= \frac{24}{36} = \frac23 .
\]
\end{example}

\section{Machten}

\begin{definition}[Macht]\label{def:g9:fractions:power}
Voor een reëel getal $a$ en een positief geheel getal $n$ is de
\emph{macht}\index{macht} $a^n$ het \omterm{def:g2:mult:def}{product} van $n$ factoren gelijk aan $a$:
\[
a^n = \underbrace{a \times a \times \dots \times a}_{n \text{ factoren}} .
\]
Per afspraak is $a^0 = 1$ (voor $a \neq 0$), en negatieve \omterm{def:g8:powers:def}{exponenten} staan voor
omgekeerden:
\[
a^{-n} = \frac{1}{a^n}.
\]
\end{definition}

\begin{example}
$2^4 = 2 \times 2 \times 2 \times 2 = 16$; \quad
$(-3)^2 = 9$, maar $-3^2 = -(3^2) = -9$ (de \omterm{def:g8:powers:def}{exponent} bindt sterker dan het
minteken); \quad $10^{-3} = \frac{1}{1000} = 0.001$; \quad
$5^1 = 5$ en $5^0 = 1$.
\end{example}

\begin{theorem}[Rekenregels voor exponenten]\label{thm:g9:fractions:rules}
Voor alle $a$, $b$ verschillend van \omterm{def:g1:counting:zero}{nul} en alle gehele getallen $m$, $n$:
\[
a^m \times a^n = a^{m+n},
\qquad
\frac{a^m}{a^n} = a^{m-n},
\qquad
\left(a^m\right)^n = a^{mn},
\qquad
(ab)^n = a^n b^n .
\]
\end{theorem}

\begin{proof}
Voor positieve \omterm{def:g8:powers:def}{exponenten} tel je de factoren. Voor de eerste regel: $a^m \times
a^n$ bestaat uit $m$ factoren $a$ gevolgd door $n$ factoren $a$, samen $m + n$
factoren. Voor de derde: $\left(a^m\right)^n$ herhaalt $n$ keer een blok van $m$
factoren, wat $mn$ factoren geeft. De twee andere regels gaan net zo; de
uitbreiding naar \omterm{def:g1:counting:zero}{nul} en naar negatieve \omterm{def:g8:powers:def}{exponenten} volgt uit
$a^{-n} = \frac{1}{a^n}$ (en maakt de afspraak $a^0 = 1$ de enige samenhangende
keuze, want $a^n \times a^0$ moet $a^{n+0} = a^n$ zijn).
\end{proof}

\begin{example}\label{ex:g9:fractions:rules}
Vereenvoudig stap voor stap:
\[
\frac{3^7 \times 3^{-2}}{3^4}
= \frac{3^{7 + (-2)}}{3^4}
= \frac{3^5}{3^4}
= 3^{5-4} = 3^1 = 3 .
\]
En met twee letters: $(2a^3)^2 \times a = 4 a^6 \times a = 4a^7$.
\end{example}

\section{Machten van tien en wetenschappelijke notatie}

\begin{proposition}[Machten van tien]\label{prop:g9:fractions:ten}
Voor een positief geheel getal $n$ schrijf je $10^n$ als ``$1$ gevolgd door $n$
nullen'', en is $10^{-n} = 0.0\dots01$ met de $1$ op de $n$-de decimale plaats.
Een \omterm{def:g6:decimals:places}{decimaal getal} met $10^n$ (respectievelijk $10^{-n}$) vermenigvuldigen
verschuift zijn \omterm{def:g4:decimals:point}{decimaalteken} $n$ plaatsen naar rechts (respectievelijk naar
links).
\end{proposition}

\begin{definition}[Wetenschappelijke notatie]\label{def:g9:fractions:scientific}
De \emph{wetenschappelijke notatie}\index{wetenschappelijke notatie} van een
positief \omterm{def:g6:decimals:places}{decimaal getal} is de enige manier om het te schrijven als
\[
a \times 10^n,
\qquad\text{met } 1 \leq a < 10 \text{ en } n \text{ een geheel getal.}
\]
\end{definition}

\begin{example}\label{ex:g9:fractions:scientific}
De afstand Aarde--Zon is ongeveer $149\,600\,000$ km $= 1.496 \times 10^8$ km. De
grootte van een watermolecule is ongeveer $0.000\,000\,000\,28$ m
$= 2.8 \times 10^{-10}$ m. Om met zulke getallen te rekenen, groepeer je de
decimale delen en de \omterm{def:g9:fractions:power}{machten} van tien:
\[
(3 \times 10^5) \times (2.5 \times 10^{-2})
= (3 \times 2.5) \times 10^{5 + (-2)}
= 7.5 \times 10^{3}.
\]
\end{example}

\begin{method}[Een getal in wetenschappelijke notatie zetten]\label{met:g9:fractions:scientific}
\begin{enumerate}
  \item Zet het \omterm{def:g4:decimals:point}{decimaalteken} juist achter het eerste cijfer dat niet \omterm{def:g1:counting:zero}{nul} is; dat
        geeft de factor $a$ met $1 \leq a < 10$;
  \item tel over hoeveel plaatsen het teken is verschoven: dat aantal is $n$,
        positief als het oorspronkelijke getal $\geq 10$ was, negatief als het
        $< 1$ was;
  \item controleer: $a \times 10^n$ moet het oorspronkelijke getal teruggeven.
\end{enumerate}
\end{method}

\section{Oefeningen}

\begin{exercise}[$\star$]\label{exo:g9:fractions:1}
Bereken en geef het resultaat als een volledig vereenvoudigde \omterm{def:g9:fractions:fraction}{breuk}:
\[
\frac13 + \frac15, \qquad
\frac72 - \frac53, \qquad
\frac{5}{12} + \frac38, \qquad
3 - \frac45 .
\]
\end{exercise}

\begin{exercise}[$\star$]\label{exo:g9:fractions:2}
Bereken en vereenvoudig:
\[
\frac59 \times \frac{27}{35}, \qquad
\frac{8}{15} \div \frac{4}{5}, \qquad
\frac23 \times \frac34 \times \frac45 .
\]
\end{exercise}

\begin{exercise}[$\star$]\label{exo:g9:fractions:3}
Bereken:
\[
2^5, \qquad (-2)^4, \qquad -2^4, \qquad 4^{-2}, \qquad
\left(\tfrac23\right)^3, \qquad 7^0 .
\]
\end{exercise}

\begin{exercise}[$\star$]\label{exo:g9:fractions:4}
Schrijf als één \omterm{def:g9:fractions:power}{macht}:
\[
5^3 \times 5^6, \qquad
\frac{7^9}{7^5}, \qquad
\left(2^3\right)^4, \qquad
\frac{3^2 \times 3^7}{3^5}, \qquad
2^5 \times 4 .
\]
\end{exercise}

\begin{exercise}[$\star$]\label{exo:g9:fractions:5}
Schrijf in \omterm{def:g9:fractions:scientific}{wetenschappelijke notatie}: $45\,000$; $0.0072$; $302.5$;
$0.000\,000\,9$; twaalf miljoen.
\end{exercise}

\begin{exercise}[$\star\star$]\label{exo:g9:fractions:6}
Bereken, en geef het resultaat in \omterm{def:g9:fractions:scientific}{wetenschappelijke notatie}:
\[
(2 \times 10^7) \times (4.5 \times 10^{-3}),
\qquad
\frac{9 \times 10^5}{3 \times 10^{-2}},
\qquad
(5 \times 10^3)^2 .
\]
\end{exercise}

\begin{exercise}[$\star\star$]\label{exo:g9:fractions:7}
Een tank is voor $\frac35$ vol. Na het toevoegen van $30$ liter is hij voor
$\frac34$ vol. Wat is de \omterm{def:g2:measure:units}{inhoud} van de tank? (Druk eerst uit welk deel van de tank
erbij is gekomen.)
\end{exercise}

\begin{exercise}[$\star\star$]\label{exo:g9:fractions:8}
Alice eet $\frac14$ van een cake op, daarna eet Ben $\frac13$ van wat overblijft,
en daarna eet Carla $\frac12$ van wat er dan nog over is. Welk deel van de cake
blijft er aan het einde over? (Bereken na elke stap de \omterm{def:g3:division:remainder}{rest}.)
\end{exercise}

\begin{exercise}[$\star\star$]\label{exo:g9:fractions:9}
Licht legt $3 \times 10^8$ meter per seconde af. De Zon staat ongeveer
$1.5 \times 10^{11}$ m ver. Hoe lang doet zonlicht erover om ons te bereiken? Geef
het antwoord in seconden (\omterm{def:g9:fractions:scientific}{wetenschappelijke notatie} is niet nodig), en daarna in
minuten en seconden.
\end{exercise}

\begin{exercise}[$\star\star\star$]\label{exo:g9:fractions:10}
Wat is groter, $2^{30}$ of $3^{20}$? (Tip: schrijf beide als \omterm{def:g9:fractions:power}{machten} met \omterm{def:g8:powers:def}{exponent}
$10$ met behulp van $\left(a^m\right)^n = a^{mn}$, en vergelijk $2^3$ met $3^2$.)
\end{exercise}

\section{Opgave: de geheime code van de repeterende decimalen}

\begin{problem}[{Weekendopgave --- elke breuk heeft een decimale schrijfwijze die
stopt of repeteert, elke repeterende decimale schrijfwijze is een breuk, en
$0.999\ldots$ is precies $1$}]
\label{pb:g9:fractions:1}
Deel $3$ door $11$ en de cijfers vallen in een dreun: $0.272727\ldots$, eindeloos.
Deel $1$ door $4$ en de cijfers stoppen abrupt: $0.25$. Deze opgave bewijst dat
dit de enige twee mogelijke gedragingen van een \omterm{def:g9:fractions:fraction}{breuk} zijn --- en zet daarna de
machine in omgekeerde richting: gegeven een repeterende decimale schrijfwijze,
bouwt ze de \omterm{def:g9:fractions:fraction}{breuk} terug die erachter schuilt. Onderweg beslecht ze, eens en voor
altijd, de meest bediscussieerde gelijkheid van de schoolwiskunde.

\medskip
\noindent\textbf{Deel I --- Van \omterm{def:g9:fractions:fraction}{breuk} naar decimale schrijfwijze.}
\begin{enumerate}
  \item Bereken met een staartdeling de decimale schrijfwijzen van $\frac14$,
        $\frac13$, $\frac56$ en $\frac{3}{11}$. Welke stoppen, welke repeteren, en
        met welk repeterend blok?
  \item Leg uit waarom de decimale schrijfwijze van \emph{elke} \omterm{def:g9:fractions:fraction}{breuk} $\frac ab$
        ofwel moet stoppen, ofwel vanaf een zeker punt moet repeteren. (Welke
        waarden kunnen de opeenvolgende \omterm{def:g3:division:remainder}{resten} in de staartdeling door $b$
        aannemen? Wat gebeurt er op het ogenblik dat een \omterm{def:g3:division:remainder}{rest} voor de tweede keer
        opduikt?)
  \item Zet de deling van $1$ door $7$ ver genoeg voort om het repeterende blok te
        vinden. Hoe lang is het, en hoe verhoudt die lengte zich tot de grens die
        vraag 2 belooft?
  \item Sommige \omterm{def:g9:fractions:fraction}{breuken} stoppen doordat hun \omterm{def:g4:fractions:def}{noemer} tot een \omterm{def:g9:fractions:power}{macht} van tien kan
        worden opgeblazen: $\frac14 = \frac{25}{100}$,
        $\frac{3}{8} = \frac{375}{1000}$. Leg het criterium uit, en ook waarom
        geen enkele vermenigvuldiger die een heel getal is, $3$, $6$ of $11$ ooit
        in $10$, $100$, $1000, \dots$ kan omzetten (op welk cijfer zou het \omterm{def:g2:mult:def}{product}
        eindigen?).
  \item Verdeel zonder te delen $\frac{9}{40}$, $\frac{5}{12}$ en $\frac{33}{50}$
        over ``stopt'' en ``repeteert'', en geef daarna de decimale schrijfwijze
        van de twee die stoppen.
\end{enumerate}

\medskip
\noindent\textbf{Deel II --- Van repeterende decimale schrijfwijze terug naar de
\omterm{def:g9:fractions:fraction}{breuk}.}
\begin{enumerate}[resume]
  \item De verschuiftruc. Zij $x = 0.7777\ldots$ Bereken $10x$, daarna
        $10x - x$, en leid $x$ als \omterm{def:g9:fractions:fraction}{breuk} af. Ga het na door te delen.
  \item Zij $x = 0.727272\ldots$ Welke \omterm{def:g9:fractions:power}{macht} van tien verschuift $x$ over precies
        één repeterend blok? Schrijf $x$ ermee als een \omterm{def:g9:fractions:fraction}{breuk} in eenvoudigste
        vorm.
  \item Een gemengd geval: $x = 0.58333\ldots$ (alleen de $3$ repeteert). Bereken
        $1000x - 100x$ en leid $x$ als \omterm{def:g9:fractions:fraction}{breuk} in eenvoudigste vorm af.
  \item De beroemde: pas de verschuiftruc toe op $x = 0.9999\ldots$ Welke \omterm{def:g9:fractions:fraction}{breuk}
        --- welk \emph{getal} --- vind je? Verzoen het resultaat met je
        aanvoelen, en denk daarbij aan de spookbuur van
        \cref{pb:g6:decimals:1} en de krimpende \omterm{def:g3:division:remainder}{rest} van
        \cref{pb:g6:fractions:1}: ligt $0.999\ldots$ ``net onder'' $1$, of is het
        een andere naam voor $1$?
  \item Zet $2.454545\ldots$ en $0.123123123\ldots$ om in \omterm{def:g9:fractions:fraction}{breuken} in
        eenvoudigste vorm.
\end{enumerate}

\medskip
\noindent\textbf{Deel III --- Het woordenboek afgemaakt.}
\begin{enumerate}[resume]
  \item Vragen 6--10 suggereren een woordenboek: een blok van $k$ cijfers dat
        vanaf het \omterm{def:g4:decimals:point}{decimaalteken} repeteert, is dat blok gedeeld door $k$ negens
        ($0.727272\ldots = \frac{72}{99}$). Voorspel met het woordenboek
        $0.037037\ldots$ als \omterm{def:g9:fractions:fraction}{breuk}, vereenvoudig --- en verwonder je: welke
        verrassend eenvoudige \omterm{def:g9:fractions:fraction}{breuk} dreunt ``$037$''? Ga het na door te delen.
  \item Combineer het woordenboek met een verschuiving: schrijf
        $0.0272727\ldots$ als een \omterm{def:g9:fractions:fraction}{breuk} in eenvoudigste vorm.
  \item Het getal $0.101001000100001\ldots$ (een $1$, dan één $0$, dan een $1$,
        dan twee nullen, dan een $1$, dan drie nullen, en zo verder) stopt nooit.
        Is het een repeterende decimale schrijfwijze? Wat zegt deel I er dan over
        --- kan het een \omterm{def:g9:fractions:fraction}{breuk} zijn? (Je hebt net je eerste bewijsbaar
        \emph{irrationale} getal ontmoet; het beroemdste wordt in
        \cref{ch:g9:sqrt} gevangen.)
  \item Hoeveel cijfers heeft $2^{30}$ ongeveer? Gebruik
        $2^{10} = 1\,024 \approx 1.024 \times 10^3$ en de rekenregels voor
        \omterm{def:g8:powers:def}{exponenten} (\cref{thm:g9:fractions:rules}) om $2^{30}$ in (benaderde)
        \omterm{def:g9:fractions:scientific}{wetenschappelijke notatie} te schrijven, en besluit. (Vergelijk
        \cref{exo:g9:fractions:10}.)
  \item Slot: schrijf $0.20262026\ldots$ als een \omterm{def:g9:fractions:fraction}{breuk}, en formuleer de les in
        twee richtingen van de hele opgave: welke decimale schrijfwijzen zijn
        \omterm{def:g9:fractions:fraction}{breuken}, en welke \omterm{def:g9:fractions:fraction}{breuken} hebben welke decimale schrijfwijze?
\end{enumerate}
\end{problem}
